\documentclass[11pt, a4paper]{article}

% --- UNIVERSAL PREAMBLE BLOCK ---
\usepackage[a4paper, top=2.5cm, bottom=2.5cm, left=2cm, right=2cm]{geometry}
\usepackage{fontspec}
\usepackage[english, bidi=basic, provide=*]{babel}

\babelprovide[import, onchar=ids fonts]{english}

% Set default/Latin font to Sans Serif in the main (rm) slot
\babelfont{rm}{Noto Sans}

\usepackage{amsmath}
\usepackage{amssymb}
\usepackage{amsthm}
\usepackage{booktabs}
\usepackage{enumitem}
\setlist[itemize]{label=-}

\usepackage[colorlinks=true, linkcolor=blue, urlcolor=blue, citecolor=blue]{hyperref}

% Theorem environments
\newtheorem{definition}{Definition}
\newtheorem{theorem}{Theorem}
\newtheorem{lemma}{Lemma}
\newtheorem{corollary}{Corollary}
\newtheorem{proposition}{Proposition}
\newtheorem{remark}{Remark}

\title{The COD Kernel Representation Framework: \\ Operator-Algebraic Reconstruction of Quantum Correlation Geometry}
\author{Omega Protocol Formalization}
\date{April 2026}

\begin{document}

\maketitle

\begin{abstract}
We present an operator-algebraic framework that reconstructs aspects of quantum correlation geometry from Completely Positive Trace-Preserving (CPTP) informational flow. By explicitly separating functional analytic derivations from physical interpretations, we show how non-commutative probability structures can be induced from families of Choi operators. We construct a positive-semidefinite Gram kernel via the Hilbert--Schmidt inner product of Choi states and use it to define a positive functional on the operator system generated by node observables. Using the Gelfand--Naimark--Segal (GNS) construction, we obtain a cyclic representation in which the kernel's imaginary part controls commutator expectation values. For bipartite settings, we recover Tsirelson-type bounds with a formal $C^*$-algebraic proof linking tensor embeddings to Grothendieck limitations. Furthermore, we prove that when the symplectic flux vanishes, the Gram representation geometrically collapses, strictly enforcing the classical CHSH bound ($\le 2$). Finally, we expand on the asymptotic-commutativity hypothesis, demonstrating that under large-time limits, the represented algebra abelianizes and strictly reduces to a classical measure space via Gelfand duality.
\end{abstract}

\section{Introduction}

The objective of this framework is to establish a mathematically rigorous pipeline from informational operations to quantum geometric structures. To avoid circularity, this text strictly delineates rigorous functional analysis (Theorem and Definition layers) from physical correspondences (Interpretation layers).

\section{The Gram Kernel Space}

Let $\mathcal{N}$ be a discrete index set of interacting nodes. For each $i \in \mathcal{N}$, let $\mathcal{E}_i$ be a Completely Positive Trace-Preserving (CPTP) map with Choi--Jamio\l{}kowski operator $J_i$.

\begin{definition}[Informational Gram Kernel]
Assume $J_i$ are Hilbert--Schmidt operators on a common Hilbert space. Define
\begin{equation}
    K(i,j) = \langle J_i, J_j \rangle_{HS} = \mathrm{Tr}(J_i^\dagger J_j),
\end{equation}
for $i,j \in \mathcal{N}$.
\end{definition}

\begin{proposition}[Canonical Hermitian Decomposition]
As a Gram kernel over a complex Hilbert space, $K$ is Hermitian ($K(i,j)=\overline{K(j,i)}$) and positive semidefinite. Hence
\begin{equation}
    K(i,j) = g_{ij} + i\,\omega_{ij},
\end{equation}
where $g_{ij}=\operatorname{Re}K(i,j)$ is symmetric and $\omega_{ij}=\operatorname{Im}K(i,j)$ is skew-symmetric.
\end{proposition}

\begin{remark}[Mathematical Constraint]
The positivity condition $K \succeq 0$ constrains the admissible pair $(g,\omega)$. In finite dimensions, this is equivalent to positivity of the complex matrix $g+i\omega$.
\end{remark}

\begin{remark}[Physical Interpretation]
We identify $g_{ij}$ as a Chain Overlap Density (COD), representing metric overlap, and $\omega_{ij}$ as a Reverse Chain Overlap Density (RCOD), representing directional causal flux.
\end{remark}

\section{Emergence of Observables via GNS Representation}

\begin{definition}[Operator-System State Induced by the Kernel]
Let $\mathcal{V}=\operatorname{span}\{1,A_i,A_i^\dagger, A_i^\dagger A_j\,:\, i,j\in\mathcal{N}\}$ be the operator system generated up to degree two. Define
\begin{equation}
\omega_{\mathrm{state}}(A_i^\dagger A_j):=K(i,j), \qquad \omega_{\mathrm{state}}(1):=1,
\end{equation}
and extend linearly on $\mathcal{V}$. Positivity on degree-two elements follows from $K\succeq 0$.
\end{definition}

\begin{theorem}[GNS Emergence]
If $\omega_{\mathrm{state}}$ admits a positive extension to a unital $C^*$-algebra $\mathcal{A}$ containing $\mathcal{V}$, then by the GNS construction there is a cyclic triple $(\pi,\mathcal{H}_{\mathrm{GNS}},\Omega)$ with
\begin{equation}
    \omega_{\mathrm{state}}(X)=\langle \Omega,\pi(X)\Omega\rangle, \quad X\in\mathcal{A},
\end{equation}
and in particular
\begin{equation}
    K(i,j)=\langle \Omega,\pi(A_i)^\dagger\pi(A_j)\Omega\rangle.
\end{equation}
\end{theorem}

\section{Symplectic Form and the Central Extension}

\begin{proposition}[Expectation-Value 2-Form]
For represented observables $\hat E_i:=\pi(A_i)$,
\begin{equation}
    \langle\Omega,[\hat E_i,\hat E_j]\Omega\rangle
    =\langle\Omega,\hat E_i\hat E_j\Omega\rangle-\langle\Omega,\hat E_j\hat E_i\Omega\rangle
    =2i\,\operatorname{Im}\langle\Omega,\hat E_i\hat E_j\Omega\rangle.
\end{equation}
If $\langle\Omega,\hat E_i\hat E_j\Omega\rangle=K(i,j)$ (e.g., for self-adjoint generators with the chosen state), then
\begin{equation}
\langle\Omega,[\hat E_i,\hat E_j]\Omega\rangle=2i\,\omega_{ij}.
\end{equation}
\end{proposition}

\begin{remark}[Operator Identity vs. Expectation Identity]
The scalar $\omega_{ij}$ is an expectation-level 2-form. It does not by itself imply an operator identity $[\hat E_i,\hat E_j]=2i\omega_{ij}I$. Such identities require additional algebraic assumptions (e.g., a CCR/Weyl representation and centrality conditions).
\end{remark}

\section{Bipartite Embedding and Correlation Bounds}

To establish the limits of non-local correlation strictly from the operator-algebraic structure of the GNS representation, we formalize the tensor embedding and provide an explicit proof of the Tsirelson bound.

\begin{theorem}[Tsirelson-Type Bound in the Bipartite Representation]
Consider a bipartite representation with observables $A_0,A_1$ on subsystem $A$ and $B_0,B_1$ on subsystem $B$. Assume these are self-adjoint, contractive operators ($\|A_x\|\le 1$, $\|B_y\|\le 1$) acting on a Hilbert space $\mathcal{H}_A \otimes \mathcal{H}_B$. For the CHSH operator
\begin{equation}
\mathcal{B} := A_0\otimes B_0 + A_0\otimes B_1 + A_1\otimes B_0 - A_1\otimes B_1,
\end{equation}
it holds that $\|\mathcal{B}\| \le 2\sqrt{2}$. Consequently, every state $\rho$ on the tensor product algebra satisfies $|\operatorname{Tr}(\rho\,\mathcal{B})| \le 2\sqrt{2}$.
\end{theorem}

\begin{proof}
By the standard property of the C*-norm, $\|\mathcal{B}\| = \sqrt{\|\mathcal{B}^2\|}$. Since $A_x$ and $B_y$ are self-adjoint operators in commuting subalgebras (enforced by the tensor product $\mathcal{H}_A \otimes \mathcal{H}_B$), we expand the square of the CHSH operator:
\begin{align*}
\mathcal{B}^2 &= (A_0 \otimes B_0 + A_0 \otimes B_1 + A_1 \otimes B_0 - A_1 \otimes B_1)^2 \\
&= A_0^2 \otimes (B_0 + B_1)^2 + A_1^2 \otimes (B_0 - B_1)^2 + [A_0, A_1] \otimes [B_0, B_1]
\end{align*}
Using the contractive assumption $\|A_x^2\| \le 1$, we can upper-bound the commuting parts:
\begin{align*}
\mathcal{B}^2 &\le I \otimes (B_0 + B_1)^2 + I \otimes (B_0 - B_1)^2 + [A_0, A_1] \otimes [B_0, B_1] \\
&= 2I \otimes (B_0^2 + B_1^2) + [A_0, A_1] \otimes [B_0, B_1] \\
&\le 4I \otimes I + [A_0, A_1] \otimes [B_0, B_1]
\end{align*}
For the commutator term, we use the norm bounds $\|[A_0, A_1]\| \le 2\|A_0\|\|A_1\| \le 2$ and $\|[B_0, B_1]\| \le 2$. Thus, the norm of the tensor product of commutators satisfies:
\begin{equation}
\| [A_0, A_1] \otimes [B_0, B_1] \| = \|[A_0, A_1]\| \cdot \|[B_0, B_1]\| \le 4
\end{equation}
Applying the triangle inequality to the operator norm:
\begin{equation}
\|\mathcal{B}^2\| \le 4 + 4 = 8 \implies \|\mathcal{B}\| \le 2\sqrt{2}.
\end{equation}
This explicitly realizes the upper constraint of Grothendieck's inequality for the complex operator space generated by the informational kernel, proving the bound strictly via the C*-algebraic embedding.
\end{proof}

\begin{remark}[Physical Interpretation]
This yields a Tsirelson bound directly from the spatial properties of the embedded bipartite representation. We derive the maximum limit of quantum correlation without invoking external quantum mechanical postulates, relying only on the operator norm limitations inherited from the GNS construction.
\end{remark}

\section{Asymptotic Abelianization and the Classical Bound}

We formally connect the emergence of classical measure spaces and local realist bounds to the vanishing of the kernel's skew-symmetric component under the action of Markovian tail chains.

\begin{proposition}[Asymptotic Commutativity and Gelfand Isomorphism]
Let the sequential composition of CPTP maps define a dynamical semigroup acting on the observable algebra, denoted $A_i^{(t)}$. We impose the explicit asymptotic commutativity hypothesis representing the collapse of informational extremity:
\begin{equation}
\lim_{t\to\infty}\|[\pi(A_i^{(t)}),\pi(A_j^{(t)})]\|=0\quad\text{for all }i,j.
\end{equation}
Let $\mathcal{A}_\infty$ denote the $C^*$-algebra generated by the asymptotic operators $A_i^{(\infty)}$. The vanishing of the commutator norm implies that $\mathcal{A}_\infty$ is a commutative $C^*$-algebra. 

By the Gelfand representation theorem, there exists an isometric *-isomorphism between the commutative unital $C^*$-algebra $\mathcal{A}_\infty$ and the algebra $C(X)$ of continuous functions on a compact Hausdorff space $X$, where $X$ is the spectrum of $\mathcal{A}_\infty$.
\end{proposition}

\begin{theorem}[Classical CHSH Bound via Abelian Collapse]
Assume the asymptotic condition $\omega_{ij}^{(\infty)} = 0$ holds, rendering $\mathcal{A}_\infty$ commutative. For any four self-adjoint, contractive observables $A_0, A_1, B_0, B_1 \in \mathcal{A}_\infty$, the CHSH expectation value evaluated by the asymptotic state functional $\omega_{\mathrm{state}}^{(\infty)}$ is strictly bounded by 2:
\begin{equation}
    |\omega_{\mathrm{state}}^{(\infty)}(A_0 B_0 + A_0 B_1 + A_1 B_0 - A_1 B_1)| \le 2.
\end{equation}
\end{theorem}

\begin{proof}
By Proposition 3, the commutative $C^*$-algebra $\mathcal{A}_\infty$ is isometrically *-isomorphic to $C(X)$ via the Gelfand transform $A \mapsto \hat{a}$. By the Riesz-Markov-Kakutani representation theorem, the positive linear functional $\omega_{\mathrm{state}}^{(\infty)}$ corresponds to integration against a unique regular Borel probability measure $\mu$ on $X$. The probability measure is strictly normalized ($\int_X d\mu = 1$) because the CPTP condition guarantees $\omega_{\mathrm{state}}(1) = 1$.

The observables map to continuous functions $\hat{a}_0, \hat{a}_1, \hat{b}_0, \hat{b}_1 \in C(X)$ with ranges in $[-1, 1]$. The CHSH functional evaluates to the integral of the polynomial:
\begin{equation}
    S(x) = \hat{a}_0(x)\hat{b}_0(x) + \hat{a}_0(x)\hat{b}_1(x) + \hat{a}_1(x)\hat{b}_0(x) - \hat{a}_1(x)\hat{b}_1(x).
\end{equation}
Factoring the integrand yields $S(x) = \hat{a}_0(x)[\hat{b}_0(x) + \hat{b}_1(x)] + \hat{a}_1(x)[\hat{b}_0(x) - \hat{b}_1(x)]$. Because $\hat{b}_i(x) \in [-1, 1]$, for any $x \in X$, if $|\hat{b}_0(x) + \hat{b}_1(x)| = 2$, then $\hat{b}_0(x) - \hat{b}_1(x) = 0$, and vice versa. Therefore, $\sup_{x \in X} |S(x)| \le 2$.

Integrating this bound against the probability measure $\mu$:
\begin{equation}
    |\omega_{\mathrm{state}}^{(\infty)}(\mathcal{B})| \le \int_X |S(x)| \, d\mu(x) \le \int_X 2 \, d\mu(x) = 2.
\end{equation}
The absence of the skew-symmetric 2-form collapses the representation into a classical $L^1(X, \mu)$ space, mathematically precluding the geometric volumes necessary to violate local realism.
\end{proof}

\begin{remark}[Geometric and Normalization Constraints]
Geometrically, the RCOD ($\omega_{ij}$) acts as an orthogonal "twist" in the Gram representation, expanding the informational parallelepiped to allow quantum violations up to $2\sqrt{2}$. When this twist vanishes, the Gram vectors are forced into a flat, locally realist projection (zero symplectic volume). Crucially, the CPTP constraint ensures that the integration measure $\mu$ remains normalized across all partial traces, preventing any artificial "informational leaks" or non-physical scaling that would permit PR-box behaviors beyond the derived bounds.
\end{remark}

\section{Conclusion}

Separating Gram-kernel data from operator-representation data yields a clean mathematical pipeline:
\[
\text{CPTP maps}\to \text{Gram kernel}\to \text{positive functional}\to \text{GNS representation}\to \text{operator correlations}.
\]
Within this chain, the imaginary kernel component strictly bounds commutator expectations. By proving the Tsirelson limit mathematically via the bipartite C*-algebraic expansion, and formally deriving the classical limit via Gelfand isomorphism, we establish a robust structural framework deriving non-local correlation and classical decoherence strictly from informational transition bounds.

\begin{thebibliography}{9}
\bibitem{GNS} Gelfand, I. \& Naimark, M. (1943). On the imbedding of normed rings into the ring of operators in Hilbert space. \textit{Recueil Math\'ematique}, 12(54).
\bibitem{QuantumGrothendieck} Pisier, G. \& Shlyakhtenko, D. (2002). Grothendieck's theorem for operator spaces. \textit{Inventiones Mathematicae}, 150(1).
\bibitem{ExtremeEvents} Papastathopoulos, I. et al. (2016). Extreme Events of Markov Chains. \textit{arXiv:1510.08920}.
\end{thebibliography}

\end{document}
